\chapter{Boxes \& Environments}
\label{ch:boxes}

\section{Example Box}

The \texttt{example} environment provides a numbered, captioned box for
worked examples:

\begin{minted}{latex}
\begin{example}[A Simple Derivation]
  \caption{Differentiating $x^2$}
  \begin{align*}
    \frac{d}{dx}\,x^2 &= 2x
  \end{align*}
\end{example}
\end{minted}

Key properties:

\begin{itemize}
  \item Auto-incrementing counter (Example~1, Example~2, \ldots).
  \item Entries appear in the List of Examples (LoE) when applicable.
  \item Breakable across pages via \texttt{tcolorbox}'s \texttt{breakable}
        library.
  \item Supports \cs{caption} for a descriptive label.
\end{itemize}

\section{Key Concept Box}

Handout-focused documents can use \texttt{keyconceptbox} to highlight
essential ideas:

\begin{minted}{latex}
\begin{keyconceptbox}
  \keyconcept{The Central Limit Theorem}
  The mean of a sufficiently large sample of independent random variables
  is approximately normally distributed.
\end{keyconceptbox}
\end{minted}

\section{Callout Box}

White papers and reports benefit from the \cs{callout} macro:

\begin{minted}{latex}
\callout{Important}{This section assumes familiarity with basic
probability theory.}
\end{minted}

The callout renders as a \texttt{tcolorbox} with the active theme's accent
color.

\section{Presentation Blocks}

Five themed block environments are available for presentations and slides:

\begin{minted}{latex}
\begin{presentationblock}{Main Idea}
  Content goes here.
\end{presentationblock}

\begin{alertblock}{Warning}
  Proceed with caution.
\end{alertblock}

\begin{exampleblock}{Illustration}
  An example goes here.
\end{exampleblock}

\begin{noteblock}{Remark}
  Additional context.
\end{noteblock}

\begin{definitionblock}{Definition}
  A \emph{widget} is defined as \ldots
\end{definitionblock}
\end{minted}

Each block is styled by the active color theme (see
Chapter~\ref{ch:color-themes}).

\section{Git Verification Box}

The \cs{GitVerificationBox} command renders a colophon-style box showing the
current Git commit SHA and verification instructions:

\begin{minted}{latex}
\GitVerificationBox
\end{minted}

Output includes the full SHA, short SHA, and a shell command for the reader
to verify the build.  See Chapter~\ref{ch:git} for full details.

\section{Todo Notes}

Conditional TODO notes are available via the \texttt{todonotes} option:

\begin{minted}{latex}
\todo{Revise this section before final submission.}
\end{minted}

Notes are only rendered when the \texttt{todonotes} class option is enabled.
They are stripped automatically in production builds.

\section{Censor Box}

The \cs{censorbox} command creates a censored region for larger redactions:

\begin{minted}{latex}
\censorbox{%
  This entire paragraph contains sensitive information
  that should not appear in the published version.
}
\end{minted}

Like \cs{todo}, censor boxes are conditional on the \texttt{censoring}
class option.  See Chapter~\ref{ch:censoring} for details.
